% Paper 1 abstract (ICES JMS / Fish and Fisheries style).
% Framing: methods paper + demonstration (not a regional case-study product).
% Locked identity: one treatment (closed-loop ICES Category 1 advice
% rule vs historical realised F) + one no-feedback benchmark (constant
% published ICES F_MSY = Ftar). Demonstration: four gadoids
% (Celtic/Irish Sea cod and whiting); OM = reference Beverton--Holt (bh3).
% Scope: perfect-information residual-replay backtest (funder deliverable).
% Do not \input this into paper.tex until that draft is retired;
% paper.tex still carries the six-stock (incl. mackerel) working-paper abstract.
%
% Title (locked; set in skeleton.tex):
%   Separating advice-rule performance from implementation and
%   productivity: a historical closed-loop backtest
% Alternative (regional; carried as \thanks in skeleton.tex):
%   A backtest of the ICES Category 1 advice rule for Celtic and
%   Irish Sea gadoids
%
\begin{abstract}
Realised stock trajectories confound fishing, environmental variability,
and the performance of advice frameworks. Retrospective analysis and
management strategy evaluation (MSE) leave a gap: neither isolates an
advice rule from implementation under the productivity that actually
occurred. We present a historical closed-loop \emph{backtest} that holds
observed productivity fixed in an age-structured operating model,
contrasts perfect closed-loop implementation of a harvest control rule
(HCR), as coded, with historical realised fishing mortality, and scores
outcomes against both operational control points and operating-model (OM)
reference points. A constant open-loop path at the published ICES
$F_{\mathrm{MSY}}$ (equal to $\Ftar$ in the advice rule) is retained
only as a no-feedback benchmark. We demonstrate the design with four ICES
Category~1 gadoids in the Celtic and Irish Seas---Celtic Sea cod, Irish
Sea whiting, Irish Sea cod, and Celtic Sea whiting---under perfect
information (true OM SSB; residual replay; zero implementation error)
with the ICES $\pm20\%$ TAC stability clause above $\Btrig$.
Across the four stocks, three outcome patterns emerge:
following the rule would have rebuilt stocks that were subject to
overfishing (realised $F$ above the target fishing mortality $\Ftar$);
the opposite can occur where realised $F$ already sat below $\Ftar$; and
the two yardsticks---ICES control points such as $\Btrig$ and OM MSY
reference points---can disagree, so that a stock scored as rebuilt
against OM MSY biomass may never clear $\Btrig$ under the rule, with the
size of the gap set by the OM stock--recruitment relationship.
\end{abstract}

\noindent\textit{Keywords:} advice; harvest control rule; backtest; maximum sustainable yield; management strategy evaluation; productivity; gadoids.
